\section{怎样掌握古今词义的变化}

很多中学生在学习文言文时，感到不好掌握的，就是不少文言实词与现代的讲法不同。的确，除了象“天、地、人、手、山、水、牛、羊、东、西、十、百”等等基本词汇外，绝大多数词的词义均有所变化，有所发展，形成了古今词义的不同，甚至上古、中古和近古的词义都有所不同。我们该怎样来掌握这些古今词义有了变化的词呢？根本的方法，当然是多学一些文言文，有意识地归纳常见文言实词的古义，比较它与现代的词义有何不同，日积月累，自然会熟能生巧，如果能再攻读一些有关文言词语的专门著作，就更好了。不过，中学同学要学的功课很多，时间有限，如果就在学习语文课本中的古文基础上，能够掌握常见文言实词的古今词义，岂不更好？何况广大社会青年忙于工作，更没有专门时间用于学习文言文，如果能在了解基本常识之后，抽空自学，做到举一反三，触类旁通，又何乐而不为呢？

我想，至少以下几点对于怎样掌握古今词义的变化是很值得注意的。

\textbf{1. 掌握文言实词以单音词为主的特点，辨析古今词义的不同。}

现代汉语中双音词很多，而古代特别是先秦的文言词语，绝大多数是单音词。因此，我们在学习文言文时，除了那些叠音词（如“萧萧”、“潸潸”等）、联绵词（如“窈窕”、“踽踽”等）和少数合成词（如“妇女”、“社稷”、“弃市”、“天下”等）之外，一般要把一个字看成为一个词，而不要把连用的两个单音词误认作一个双音词。例如\ [\ 按：凡“第二部分”所举例句没有全部注明书名和作者的，均见于中学语文课本\ ]：
\begin{entry}
\item[国家：]“天下国家，本同一理。”（《方腊起义》）“国”相当现代的“国家”，“家”相当现代的“家庭”。在现代汉语中，“国家”是个合成词，不能拆开讲。

\item[地方：]“荆之地方五千里，宋之地方五百里。”（《公输》） “地”相当现代的“土地”，“方”相当现代的“方圆”，“范围”。在现代汉语中，“地方”是个合成词，不能拆开讲，意思是指某一地区或指各级行政区域。

\item[指示：]“璧有瑕，请指示王。”（《廉颇蔺相如列传》）\\
 “指”在这儿是动词，相当现代的“用手指”，\\
 “示”相当现代的“把事告诉人”，这儿有\\
 “给······看”的意思。在现代汉语中，“指示”是个合成词，不能拆开讲，意思是上级对下级或者长辈对晚辈说明处理某项问题的原则和方法，有时是指示下级或者晚辈的话或文字。

\item[妻子：]“父母不相见，兄弟妻子离散。”（《孟子·二章·庄暴见孟子》） “妻”指“老婆”，“子”指“子女”。现代的“妻子”指男子的爱人，是一个合成词，不能拆开讲。

\item[颜色：]“微察公子，公子颜色愈和。”（《信陵君窃符救赵》）古代“颜”指“额”，“色”指“颜色”。所谓“颜色”，指脸上的气色。现在的“颜色”指色彩，是个合成词。

\item[衣裳：]“初闻涕泪满衣裳。”（《闻官军收河南河北》）古代上衣为“衣”，下衣为“裳”。现在的“衣裳”统指衣服，是个合成词。
\end{entry}

这样一类古今词义容易混淆的词语相当多，我们需要留心辨析。例如：

\begin{entry}
\item[寻常：]古代八尺叫“寻”；两“寻”叫“常”。

\item[规矩：]古代“正圆之器”叫“规”，犹如现代的“圆规”；“正方之器”叫“尺”，犹如现代的“曲尺”。

\item[影响：]古代的“影”指“影子”，“响”指“回声”。

\item[消息：]汉代以前，指“消灭”和“生长”，三国以后才成为一个合成词，当“音信”讲，一直沿用到今。

\item[宣言：]古代的“宣”指“公开说出”；“言”指“话”。

\item[可以：]古代的“可”指“可以”；“以”指“用、用来”。

\item[无论：]古代的“无”指“不要”；“论”指“说”。

\item[其实：]古代的“其”是代词，指“它的”；“实”指果实，引伸为实际情况。
\end{entry}

这一类在古代汉语中经常连用的意义相关的单音词，要把它们分开解释，而不能轻率地用现代的合成词词义来解释出现在文言文中的这些词。

\textbf{2. 注意掌握某些词义明显扩大的常见文言实词。}

词义的扩大就是词义所指事物范围的扩大。古今词义扩大的常见实词，它们的古今词义有共同之点，那就是属于同一类事物，只不过现代词义的范围比古代词义扩大了罢了。例如：
\begin{entry}
\item[江：]“大江东去，浪淘尽千古风流人物。”（《念奴娇·赤壁怀古》） “江”的古义专指长江，现在泛指一般的江。

\item[河：]“将军战河北，臣战河南”（《鸿门宴》） “河”的古义专指黄河，现在泛指一般的河。

\item[美：]“芳草鲜美，落英缤纷。”（《桃花源记》） “美”的古义是“味道好”，现代扩大到形容其他事物的美好，如“美景”、“美谈”、“美感”等。

\item[收：]“（衡）阴知奸党名姓，一时收禽。”（《张衡传》） “收”的古义是“捕取”，后来扩大到一般的收取。

\item[再：]“一鼓作气，再而衰，三而竭。”（《曹刿论战》） “再”的古义是“两次”或“第二次”。现代扩大到泛指行为动作的重复。

\item[好：]“是女子不好，烦大巫妪为入报河伯，得更求好女，后日送之。”（《西门豹治邺》） “好”古义指女子貌美，现代泛指与“坏”相对的意思。

\item[中国：]“若能以吴越之众与中国抗衡，不如早与之绝。”（《赤壁之战》） “中国”原指黄河流域的中原地区，后来扩大到指祖国的全部领土，现指中华人民共和国。

\item[睡：]“童子莫对，垂头而睡。” （欧阳修《秋声赋》）“睡”的古义专指坐着打瞌睡，现代泛指“睡觉”。
\end{entry}

其它如“布”（古义指麻织品，后来扩大为棉麻等织物的通称）、“苦”（古义指“大苦菜”，后来扩大为一切苦味）、“雌、雄”（本指鸟类的阴性、阳性的区别，后来扩大到兽类，甚至扩大到自然界其他事物）、“园”（本指“果园”，后来扩大到种菜、种花、观赏动物的地方也叫园），等等，在文言文中，我们都得根据其时代的不同和前后的文意，来区别它们的不同词义。

\textbf{3. 掌握某些词义明显缩小的常见文言实词。}

词义的缩小就是词义所指事物范围的缩小。古今词义缩小的常见实词，它们的古今词义有共同之处，那就是属于同一类事物，只不过现代词义的范围比古代词义缩小罢了。例如：

\begin{entry}
\item[子：]“齐侯之子，卫侯之妻。”（《诗经·卫风·硕人》）原义指儿女，不分性别，女儿也称“子”；现代专指男孩子。

\item[谷：]“其始播百谷。”（《诗经·豳风·七月》） “谷”原义指百谷的总称，后来专指稻实。

\item[臭：]“色恶不食，臭恶不食。”（《论语·乡党》） “臭”本包括一切气味，不分好坏；现代专指恶气为臭。

\item[丈夫：]“丈夫亦爱怜其少子乎？”（《触龙说赵太后》）
“丈夫”的古义通指成年男子，现代专指对“妻子”而言的男人。

\item[兄弟：]“父子俱在军中，父归，兄弟俱在军中，兄归，独子无兄弟，归养。”（《信陵君窃符救赵》）
“兄弟”本包括“兄”和“弟”两方面，现代已演变成偏义复词，专称弟弟。

\item[丈人：]“丈人知其一，未知其二。请诉之，愿丈人垂听。”（《中山狼传》） “丈人”古义是“长老”的通称。《六书正读》：“丈，借为杖行之杖，老人持杖，故曰丈人。”到了汉代，“丈人”一词有了三种意义，一称妻之父，一称夫之母，一称妻之夫。现代专指妻之父。

\item[学者：]“古之学者必有师。”（《师说》） “学者”原指“从事学习的人”，现代专指“在学术上有一定成就的人”。

\item[亲戚：]“臣所以去亲戚而事君者，徒慕君之高义也。”（《廉颇蔺相如列传》） “亲戚”古义是指“族内亲属和族外亲属”，现代专指“族外亲属”。

\item[处分：]“处分适见意，那得自任专。”（《孔雀东南飞》） “处分”古义指一般的处理、安排，现代专指对犯错误者的处理。
\end{entry}

因此，我们在解释文言文中的词义时，既要弄清古今词义的不同，不能以今义代古文；同时也要注意某些词的意义在不同的历史时期也往往有所不同。我们一定要联系上下文意，注意区分，否则就会误解文义。

\textbf{4. 特别注意掌握一批词义明显转移的常见文言实词。}

词义的转移，就是概念的转移，也就是说一个词在古代表示某一概念，后来转移到表示另一概念。这种词义转移的现象，给我们学习文言文带来了困难，因为词义转移了，我们如果还以今义去解释古义，就一定要发生错误。这是学习文言文的一大难点。但是，我们在中学语文课本中已经学习了百把篇古代诗文，接触到的文言实词上千个，如果我们在学习过程中，有意识地总结常见文言实词的词义转移的规律，并且做一些灵活运用这些规律的练习，自可达到举一反三、读懂一般简易古文的目的。例如：

\begin{entry}
\item[去：]“逝（誓）将去女（汝），适彼乐土。”（《硕鼠》）
“自去吏职，五载复还。”（《张衡传》） “因擢其金而去，吏捕得之。”（《郑人攫金》） 我们从中可以看出，“去”的古义是“离开”，这跟现代所说的“从所在的地方到别的地方”的“去”，意义刚好相反。

\item[走：]“兔走触株，折颈而死。”（《守株待兔》） “永之人争奔走焉。”（《捕蛇者说》） “与樊哙、夏侯婴、靳强、纪信等四人持剑盾步走。”（《鸿门宴》） 我们从中可以看出，“走”的古义是“跑”，这跟现代所说的“人或鸟兽的脚交互向前移动”的意思，是大不相同的。

\item[涕：]“今当远离，临表涕零，不知所言。”（《出师表》） “百姓闻之，知与不知，无老壮皆为之垂涕。”（《史记·李将军列传》） “儿涕而去。”（《促织》） “涕”的古义是“眼泪”，作动词讲是“流泪”的意思。这与现代所说的“鼻液”是大不相同的。

\item[汤：]“农夫心内如汤煮，公子王孙把扇摇。”（《智取生辰纲》） “至舍，四支（肢）僵劲不能动，腰人持汤沃灌……” （《送东阳马生序》） “那跟来的人讨了汤桶，自行烫酒。” （《林教头风雪山神庙》） 我们从中可以看出，“汤”的古义是“热水”，这跟现代所说的“羹汤”是大不相同的。

\item[卑鄙：]“先帝不以臣卑鄙，猥自枉屈，三顾臣于草庐之中······” （《出师表》） “卑鄙”的古义是“低微而鄙俗”的意思，这与现代所说的“品质恶劣”是大不相同的。

\item[牺牲：]“牺牲玉帛，弗敢加也，必以信。” （《曹刿论战》） “牺牲”古指祭祀时所用的祭品牛、羊、豕，后来词义转移为“捐弃”，现在则指“为正义事业而献出生命”，词义是大不相同的。
\end{entry}

其它如“股”（古义“大腿”，今指“臀部”）、脚（古义“小腿”，今指“足部”）、“替”（古义“衰微”，今指“代替”）、“快”（古义“痛快”，今指“迅速”）、“书”（古义“写”或者“书信”，今指“装订成册的著作”）、“让”（古义“用语言责备对方”，今指“谦让”）、“购”（古义“重金收买、重赏征求”）、“购”的对象往往不是商品；今指“买”，“买”的对象一般是商品、“对”（古义“回答”，今指“正确”）、“更衣”（古指“上厕所”的婉词，现指“换衣服”）、“交通”（古义“互相通达”，现指来往或运输）、“故事”（古义“历史上的旧事、先例”，现指真实的或虚构的用作讲述对象的事情，属于一种文学样式）、“少年”（古义指未满二十岁的青年，今指十岁左右到十五六岁的人生阶段）、“结束”（古义指“整装”，今义指“过程完结”）、“可怜”（古义“可爱”，今义“令人怜悯”）、“束手”（古义“投降”，今义“没办法”）、“市井”（古指“街市”或“商品交易的场所”，今已不用此词）、“鼓刀”（古义“动力宰杀牲畜”，今已不用此词）、“字下”（古义“房檐下面”，今已不用此词）、“奉”（古义“恭敬地拿着”，现代已不单用，仅用作词素），等等，不一而足。我们在学习文言文时，只要注意积累和比较、归纳，词义转移的常见文言实词终究是可以掌握得住的。
\newpage